% ID: SEL-D0020-OB-E01
% ESTADO: validada. S021, revisada con los criterios oficiales.
\textbf{a)} Llamemos $M$ y $N$ a las matrices de la derecha en ese
orden. Restando la primera ecuación de cinco veces la segunda,
obtenemos $13A=5N-M$. Después $B=3A-N$. Así,
\[
\boxed{A=\begin{pmatrix}1&1\\1&-1\end{pmatrix}},\qquad
\boxed{B=\begin{pmatrix}-1&0\\-1&-2\end{pmatrix}}.
\]
Al sustituir, $2A-5B=\begin{pmatrix}7&2\\7&8\end{pmatrix}$ y
$3A-B=\begin{pmatrix}4&3\\4&-1\end{pmatrix}$.

\textbf{b)} Como $\det C=5$, despejamos multiplicando \emph{por la derecha}
por $C^{-1}$, para que $CC^{-1}=I_2$. Resulta
\[
X=(I_2+D^2)C^{-1}
=\frac15\begin{pmatrix}0&2\\-2&4\end{pmatrix}
\begin{pmatrix}1&2\\-1&3\end{pmatrix}
=\boxed{\frac15\begin{pmatrix}-2&6\\-6&8\end{pmatrix}}.
\]
